\scanpage{147}
\uline{Proposition} $(E,q)$ forme quadratique régulière sur schéma $S$, \inserted{$\varphi$ forme bil. associée,} $F$ sous-\deleted{fibré} \inserted{Module strict} tot. \uline{isotrope} de $E$.

a) L'hom. composé
\[
E\xrightarrow[\substack{\text{défini}\\\text{par }\varphi}]{\Psi}E^\vee\xrightarrow{\mathrm{restr}}F^\vee
\]
est épi, donc son noyau $F^\perp$ est un sous-\deleted{fibré} Module strict de $E$, dont la formation commute à tout chgt de base.

b) \inserted{Soit $G=E/F^\perp$.} La forme $\varphi$ induit un accouplement
\[
F\times G\longrightarrow\underline{\mathcal O}_S
\]
qui est une dualité, et induit des isom.
\[
G\xrightarrow{\sim}\check F,\quad F\xrightarrow{\sim}\check G
\]

\deleted{c) Soit $\underline{\mathrm{Split}}(E,F^\perp)$ \inserted{ou simplement $\underline{\mathrm{Split}}$,} le faisceau des scindages de la suite exacte}
\diagram{\begin{tikzcd}[ampersand replacement=\&,row sep=tiny,column sep=small]
\deleted{0}\arrow[r]\&\deleted{F^\perp}\arrow[r,dash]\&\deleted{E}\arrow[r]\&\deleted{G}\arrow[r]\&\deleted{0},\\
\&\deleted{\cup}\\[-.8em]
\&\deleted{F}
\end{tikzcd}}
\deleted{qui est un torseur sous $\underline{\mathrm{Hom}}(G,F^\perp)$, soit $\underline{\mathrm{Split}}\,\underline{\mathrm{isot}}(E,F^\perp)$ \inserted{ou $\underline{\mathrm{Splitisot}}$,} le sous-faisceau des scindages \deleted{\unsure{$\sigma$}} \uline{isotropes}, i.e. tels que $\theta(G)$ soit \deleted{tot.} isotrope dans $E$, \deleted{\unsure{Cons}}}
\[
\deleted{\underline{\mathrm{Splitisot}}(E,F^\perp)\lhook\joinrel\longrightarrow\underline{\mathrm{Split}}(E,F^\perp)}.
\]
\deleted{Considérons d'autre part la suite exacte}
\[
\deleted{0\longrightarrow\underline{\mathrm{Hom}}(G,F)\longrightarrow\underline{\mathrm{Hom}}(G,F^\perp)\longrightarrow\underline{\mathrm{Hom}}(G,Q)\longrightarrow0}
\]
\deleted{où}
\[
\deleted{Q=F^\perp/F}.
\]
% archive sheet number: 146
